% !TeX root = ../collection-solutions.tex

\titledquestion{Should we tax capital income or wealth?}

\citet{boar2023income-or-wealth} ask which asset-side tax instrument a welfare-maximizing planner should use: a tax on capital income or a tax on wealth.

\begin{parts}
    \part[4] What is the key takeaway of the paper? Explain how the authors arrive at this conclusion: what kind of model do they use, which features of the US data does the calibration target, and why does this calibration matter for the result?

    \begin{solutionorbox}[12cm]
        \emph{Takeaway} (2 points): Taxing capital income beats taxing wealth. A utilitarian planner restricted to one instrument achieves a consumption-equivalent welfare gain of $\approx 9.5\%$ with a (high) capital income tax but only $\approx 5.6\%$ with a wealth tax — and when both instruments are available, the optimal wealth tax is \emph{zero}. The efficiency gains from wealth taxation are real but small, and are swamped by the redistributive power of taxing entrepreneurs' profits.

        \emph{How they get there} (2 points):
        \begin{itemize}
            \item A general-equilibrium model with entrepreneurs facing collateral constraints (as in Guvenen et al.) — but alongside an \emph{unconstrained corporate sector}, and with entrepreneurial technology that uses labor and has decreasing returns.
            \item The calibration targets how much of US wealth and income private business owners actually hold (SCF: entrepreneurs are $\approx 12\%$ of households but hold $44\%$ of wealth and $28\%$ of income). In the calibrated economy, misallocation losses are only $\approx 1.3\%$ of output — capital that constrained entrepreneurs cannot absorb flows to the unconstrained corporate sector, and persistently productive entrepreneurs save their way out of their constraints.
            \item Optimal-policy experiment: one-time permanent tax reforms, with welfare computed \emph{including the transition} and revenue rebated as lump-sum transfers. Because the efficiency pie is tiny while the redistribution stakes are large, the instrument that targets rich entrepreneurs' profits (the capital income tax) wins.
        \end{itemize}
        Grading: 1 point for the ranking (capital income tax dominates, optimal wealth tax $\approx 0$), 1 point for a headline magnitude; 1 point for the model/calibration (corporate sector, entrepreneurs' wealth and income shares), 1 point for connecting small misallocation to the result.
    \end{solutionorbox}

    \part[4] Explain intuitively why, in this model, a capital income tax redistributes from entrepreneurs to workers while a flat wealth tax does not — even though both are taxes ``on the rich'' — and why the efficiency cost of the capital income tax is small here.

    \begin{solutionorbox}[12cm]
        \begin{itemize}
            \item (1 point) Capital income in this economy includes the \emph{profits} of private business owners, who dominate the top of the wealth and income distributions. A capital income tax therefore automatically concentrates its burden on rich, high-return entrepreneurs: per unit of wealth, an entrepreneur earning excess profits pays more than a bondholder.
            \item (1 point) A flat wealth tax taxes productive and unproductive wealth holders identically. It does improve the \emph{allocation} of capital (it shifts the burden off productive entrepreneurs — the use-it-or-lose-it channel), but it forgoes precisely the redistribution from rich entrepreneurs to workers that the planner values.
            \item (1 point) The efficiency cost of taxing capital income is small because misallocation is small to begin with ($\approx 1.3\%$ of output): the unconstrained corporate sector absorbs any capital the entrepreneurial sector cannot use, and persistent ability lets productive entrepreneurs eventually self-finance.
            \item (1 point) Since the efficiency pie is about $1\%$ of output while redistribution moves welfare by an order of magnitude more (the bottom third gains $\approx 30\%$ under the capital income tax, financed by large lump-sum transfers), the capital income tax dominates. The deeper lesson: wealth versus capital income taxation cannot be ranked on efficiency grounds alone — distributional incidence is quantitatively decisive.
        \end{itemize}
    \end{solutionorbox}
\end{parts}
